\documentclass[11pt]{article}
\usepackage[margin=1in]{geometry}
\usepackage{amsmath,amssymb}
\usepackage[hidelinks]{hyperref}

\title{Disproof of Conjecture 00000000905\\
\large A non-cyclotomic second-adjacency polynomial}
\author{}
\date{}

\begin{document}
\maketitle

The principal graph \(D_4\) gives a counterexample. Both the
adjacency square and the adjacency matrix for vertices at exact
distance two have an eigenvalue that is not a root of unity.
This also holds on their even parity blocks.

\section*{A subfactor realizing \(D_4\)}

Take an outer action of \(C_3\) on the hyperfinite
\(\mathrm{II}_1\) factor \(R\), and set \(M=R\rtimes C_3\).
The index-three inclusion \(R\subset M\) is the \(D_4\) case
of the group subfactor construction \cite[Sections~1 and~2.3]{izumi}.
In bimodule terms, the basic \(R\)-\(M\) sector
\(\kappa={}_R L^2(M)_M\) is irreducible because the action is outer.
The \(R\)-\(R\) bimodule \(\kappa\overline{\kappa}\) decomposes
by the crossed-product Fourier expansion into the three
automorphism bimodules, one for each element of \(C_3\).
These are distinct irreducibles by outerness. Tensoring any of
them with \(\kappa\) gives \(\kappa\), since the corresponding
automorphism is implemented by its group unitary in \(M\).
Thus there is one odd vertex, connected once to each of three
even vertices. This is \(D_4\) with an even leaf as root,
and the subfactor has finite depth.

\section*{The second adjacency calculations}

Order the center before the three leaves. Let \(J_3\) denote the
\(3\times3\) all-ones matrix and \(\mathbf 1\) the column of three
ones. The adjacency and its square are
\[
  A=\begin{pmatrix}0&\mathbf 1^{\mathsf T}\\
                   \mathbf 1&0_{3\times3}\end{pmatrix},
  \qquad A^2=[3]\oplus J_3 .
\]
In particular, the even block of \(A^2\) is the product of the
incidence column with its transpose. On the other hand, distinct
vertices at distance two are exactly pairs of leaves. Their
adjacency is \(D_2=[0]\oplus(J_3-I_3)\). The off-diagonal
two-step-walk counts give this same matrix.

The matrix \(J_3\) has eigenvalue \(3\) on \(\mathbf 1\) and
vanishes on the two-dimensional subspace of coordinate sum zero.
On these subspaces \(J_3-I_3\) acts by \(2\) and \(-1\),
respectively. Consequently the characteristic polynomials are
\[
\begin{array}{c|cc}
 & \text{Full matrix} & \text{Even block}\\ \hline
A^2 & t^2(t-3)^2 & t^2(t-3)\\
D_2 & t(t-2)(t+1)^2 & (t-2)(t+1)^2 .
\end{array}
\]
In the first row \(3\) is a root, and in the second row \(2\)
is a root. A positive power of either integer is greater than
one, while every root of a cyclotomic polynomial is a root
of unity. Therefore none of the four characteristic polynomials
is a product of cyclotomic polynomials. The odd blocks
\([3]\) and \([0]\) also have roots that are not roots of unity.

\section*{Lean verification}

The Lean~4.19.0 project constructs \(A\) from the star adjacency
relation. It checks the two block formulas, the displayed
eigenvectors and the non-root-of-unity powers. It computes the
four characteristic polynomials by Laplace expansion with
integer coefficient lists and evaluates them at the indicated roots.
Run \verb|lake build| in \verb|lean/|
to verify the finite calculations.

\begin{thebibliography}{1}
\bibitem{izumi}
M.~Izumi, \emph{Group-subgroup subfactors revisited},
Quantum Topology \textbf{15} (2024), 721--778.
\href{https://doi.org/10.4171/QT/211}{DOI}.
\end{thebibliography}

\end{document}
